% ============================================================================
%  B/08-doc-giua-cac-dong — file chương
%  Ngày tạo: 2026-09-06 | Sửa gần nhất: 2026-09-06 | Ôn gần nhất: chưa
%  Dependency: A/01, B/06. Mức độ: cốt lõi — chương quan trọng nhất phần B.
% ============================================================================
\documentclass[../../english.tex]{subfiles}
\begin{document}
\chapter{Đọc giữa các dòng: hàm ý, dè dặt và thái độ (Reading Between the Lines)}
\ptrtarget{B/08}\ptrtarget{B/08-doc-giua-cac-dong}
\dependency{\ptr{A/01} cho tầng thứ ba của nghĩa --- ý định người viết; \ptr{B/06} cho bộ khung lập luận. Đây là chương trung tâm cho mục tiêu ``hiểu đúng ý người nói''.}

\begin{tomtat}
Trong văn học thuật tiếng Anh, phần lớn thông tin về \emph{mức tin cậy} và \emph{thái độ} không nằm ở nội dung câu mà nằm ở cách chọn từ. Chương này giải mã tầng đó. Nội dung gồm: nguyên lý hợp tác của Grice và cách suy ra điều người viết ngụ ý mà không nói; bảng đầy đủ các phương tiện dè dặt và các phương tiện nhấn mạnh, xếp theo thang cam kết; các quy ước phê bình lịch sự trong học thuật cùng bảng giải mã ``họ viết thế này thường nghĩa là gì''; cách đọc nhận xét phản biện; và ý nghĩa của những gì \emph{không} được nói. Nếu chỉ nhớ một điều: khi tác giả viết dè dặt hơn mức bằng chứng cho phép, đó là lịch sự; khi họ viết dè dặt \emph{đúng} mức, đó là thông tin --- và học đọc khác biệt giữa hai điều đó là kỹ năng đọc học thuật quan trọng nhất.
\end{tomtat}

\begin{nentang}
\kn{hàm ý (implicature)}{điều người nói ngụ ý nhưng không nói ra, suy được từ việc họ chọn nói cái này mà không nói cái kia.}
\kn{nguyên lý hợp tác}{giả định ngầm rằng người nói đang cố gắng giao tiếp có ích: nói đủ, nói thật, nói liên quan, nói rõ. Khi họ có vẻ vi phạm, người nghe suy ra hàm ý.}
\kn{dè dặt (hedging)}{làm nhẹ mức cam kết: \emph{may}, \emph{suggest}, \emph{appear to}, \emph{relatively}, \emph{in our experiments}.}
\kn{nhấn mạnh (boosting)}{làm mạnh mức cam kết: \emph{clearly}, \emph{demonstrate}, \emph{establish}, \emph{must}.}
\kn{dấu hiệu thái độ}{từ bộc lộ đánh giá của tác giả: \emph{surprisingly}, \emph{unfortunately}, \emph{importantly}, \emph{remarkably}.}
\kn{nói giảm (understatement)}{phát biểu nhẹ hơn thực tế để tạo hiệu ứng, phổ biến trong văn học thuật Anh: \emph{not uncommon}, \emph{somewhat surprising}.}
\end{nentang}

\begin{khoigoi}
\h{Một bài viết \emph{our method is fast} trong phần tóm tắt và không nhắc gì tới độ chính xác. Bạn suy ra được gì, và dựa trên nguyên tắc nào?}
\h{Vì sao văn khoa học lại đầy những từ dè dặt, trong khi lẽ ra khoa học phải nói thẳng?}
\h{\emph{Further work is needed to\ldots} là câu quen thuộc trong mọi bài báo. Nó thực sự nghĩa là gì?}
\h{Một phản biện viết \emph{the paper would benefit from a comparison with X}. Đó là gợi ý hay yêu cầu?}
\h{Có những thông tin nằm ở chỗ tác giả \emph{không} nói. Làm sao nhận ra chúng?}
\end{khoigoi}

Hãy bắt đầu bằng một ví dụ ngắn. Trong phần công trình liên quan, bạn đọc:

\begin{quote}
\emph{Smith et al.\ report impressive results on indoor sequences.}
\end{quote}

Câu này khen hay chê? Về mặt chữ, nó khen. Nhưng người đọc có kinh nghiệm đọc ra một điều khác: tác giả vừa \emph{giới hạn} kết quả đó vào ``chuỗi dữ liệu trong nhà''. Nếu phương pháp kia cũng chạy tốt ngoài trời thì đã không có chữ \emph{indoor}. Và vì bài đang viết có lẽ làm việc ngoài trời, câu này chuẩn bị cho việc chỉ ra khoảng trống ở câu sau.

Không có từ nào trong câu mang nghĩa tiêu cực. Thông tin nằm ở \emph{lựa chọn nói cái gì} --- và ở việc không nói phần còn lại. Chương này dạy cách đọc tầng đó một cách có hệ thống.

\section{Nguyên lý hợp tác: vì sao im lặng cũng có nghĩa}

Nhà triết học ngôn ngữ Paul Grice đưa ra một nhận xét đơn giản mà có hệ quả rất xa: khi hai người giao tiếp, cả hai ngầm giả định rằng người kia đang \emph{hợp tác} --- nói đủ lượng thông tin cần thiết, nói điều mình tin là đúng và có bằng chứng, nói điều liên quan, và nói cho rõ ràng\cite{grice1975}. Chính vì có giả định đó mà người nghe suy ra được những điều không được nói.

Cơ chế: nếu người nói \emph{có thể} nói một điều mạnh hơn mà lại không nói, người nghe suy ra rằng họ không nói được. Đó là toàn bộ lời giải cho câu hỏi mở đầu thứ nhất. Khi một bài viết \emph{our method is fast} trong tóm tắt --- nơi tác giả có mọi động cơ để nêu ưu điểm --- và không nhắc độ chính xác, người đọc suy ra rằng độ chính xác không phải điểm mạnh. Câu văn không nói dối; nó chỉ im lặng ở đúng chỗ.

Hệ quả này áp dụng ở khắp nơi trong văn học thuật, và nó cho một thói quen đọc rất hiệu quả: \emph{với mỗi khẳng định, hãy hỏi tác giả đã có thể nói gì mạnh hơn mà họ không nói}. Ví dụ: \emph{We evaluate on three datasets} --- vì sao ba mà không phải bộ chuẩn có bảy? \emph{The method converges in most cases} --- ``phần lớn'' là bao nhiêu, và các trường hợp còn lại thì sao?

\section{Thang cam kết đầy đủ}

Chương \ptr{A/01} đã giới thiệu các động từ báo cáo. Đây là bảng đầy đủ hơn, gồm cả các phương tiện làm nhẹ và làm mạnh, để bạn định vị bất kỳ câu nào trên một thang.

\begin{center}\footnotesize
\rowcolors{2}{white}{tblalt}
\begin{longtable}{@{}L{3.2cm}L{5.0cm}L{6.2cm}@{}}
\toprule
\textbf{Loại phương tiện} & \textbf{Ví dụ} & \textbf{Đọc ra điều gì}\\
\midrule
\endfirsthead
\toprule
\textbf{Loại phương tiện} & \textbf{Ví dụ} & \textbf{Đọc ra điều gì (tiếp)}\\
\midrule
\endhead
Động từ tình thái & \emph{may, might, could, can, should, must, will} & Mức chắc chắn; \emph{could} nhẹ nhất, \emph{must} mạnh nhất (\ptr{A/03})\\
Động từ nhận thức & \emph{suggest, indicate, appear, seem, tend to} & Tác giả không cam kết hoàn toàn; dữ liệu chưa dứt khoát\\
Trạng từ dè dặt & \emph{possibly, arguably, apparently, presumably, relatively} & \emph{Arguably} báo hiệu điểm tranh cãi; \emph{apparently} báo thông tin gián tiếp\\
Từ xấp xỉ & \emph{approximately, roughly, about, some, most, several} & Cẩn thận với \emph{most} và \emph{several}: chúng thay cho con số mà tác giả chọn không nêu\\
Giới hạn phạm vi & \emph{in our experiments, under these assumptions, for the datasets considered} & Rào chắn rõ ràng: đừng suy rộng ra ngoài\\
Quy nguồn & \emph{according to X, as reported by Y} & Tác giả \emph{không} bảo lãnh cho khẳng định này\\
Nhấn mạnh & \emph{clearly, obviously, undoubtedly, indeed} & Cẩn thận: \emph{clearly} đôi khi thay cho lập luận còn thiếu\\
Động từ mạnh & \emph{demonstrate, establish, prove, confirm} & Cam kết cao; nếu bằng chứng không tương xứng, đó là điểm yếu\\
Dấu hiệu thái độ & \emph{surprisingly, unfortunately, importantly, remarkably} & Đánh giá của tác giả, không phải dữ liệu; \emph{surprisingly} ngụ ý kết quả trái kỳ vọng của cộng đồng\\
Nói giảm & \emph{not uncommon, not unreasonable, somewhat} & Thường mạnh hơn nghĩa đen; \emph{not uncommon} nghĩa là ``khá phổ biến''\\
\bottomrule
\end{longtable}
\end{center}

Câu hỏi mở đầu thứ hai hỏi vì sao khoa học lại đầy dè dặt. Có ba lý do và cả ba đều chính đáng. \emph{Lý do nhận thức:} kết quả thực nghiệm luôn có phạm vi hữu hạn, nên viết dè dặt là mô tả đúng mức bất định. \emph{Lý do xã hội:} khẳng định mạnh mà sai làm hỏng uy tín, còn khẳng định dè dặt thì vẫn đứng vững nếu sau này có phản ví dụ. \emph{Lý do tương tác:} dè dặt để lại chỗ cho người đọc không đồng ý mà không phải bác bỏ tác giả --- một cơ chế lịch sự trong tranh luận\cite{hyland1998}.

Hệ quả cho việc đọc rất quan trọng và tinh tế: bạn cần phân biệt \emph{dè dặt đúng mức} --- mang thông tin về giới hạn thật của kết quả --- với \emph{dè dặt theo lệ} --- chỉ là quy ước thể loại. Cách phân biệt: xem dè dặt có \emph{cụ thể} không. \emph{In our experiments} là dè dặt có nội dung; \emph{it may be argued that} thường chỉ là nghi thức.

\section{Bảng giải mã phê bình lịch sự}

Trong học thuật tiếng Anh, phê bình trực tiếp bị coi là thiếu chuyên nghiệp, nên nó được mã hoá. Bảng dưới là những mẫu hay gặp nhất. Cần nói rõ: đây là \emph{quy ước diễn giải}, không phải quy tắc cứng --- ngữ cảnh vẫn quyết định, và mức độ gay gắt thay đổi theo ngành và theo tác giả.

\begin{center}\footnotesize
\rowcolors{2}{white}{tblalt}
\begin{longtable}{@{}L{6.4cm}L{7.8cm}@{}}
\toprule
\textbf{Họ viết} & \textbf{Thường có nghĩa là}\\
\midrule
\endfirsthead
\toprule
\textbf{Họ viết} & \textbf{Thường có nghĩa là (tiếp)}\\
\midrule
\endhead
\emph{X has received relatively little attention} & Chưa ai làm --- và đó là khoảng trống bài này lấp\\
\emph{Further work is needed to\ldots} & Vấn đề chưa được giải quyết; đôi khi: chúng tôi không giải quyết được\\
\emph{These results are somewhat inconsistent with\ldots} & Chúng mâu thuẫn với kết quả trước\\
\emph{The paper would benefit from\ldots} & Hãy làm điều này; trong phản biện, đây thường là \emph{yêu cầu}\\
\emph{It is not entirely clear how\ldots} & Chỗ này viết không rõ, hoặc lập luận có lỗ hổng\\
\emph{We note that method X assumes\ldots} & Phê bình: giả định đó không luôn đúng\\
\emph{While interesting, the approach\ldots} & Chuẩn bị cho một phê bình nghiêm túc\\
\emph{The comparison is limited to\ldots} & So sánh không đầy đủ\\
\emph{Results are reported on a single sequence} & Bằng chứng quá mỏng để kết luận\\
\emph{Performance is competitive with\ldots} & Không tốt hơn, chỉ ngang\\
\emph{We were unable to reproduce\ldots} & Nghi ngờ nghiêm trọng về kết quả gốc\\
\emph{This is beyond the scope of the present work} & Chúng tôi biết vấn đề này nhưng không xử lý\\
\emph{A promising direction for future work} & Chúng tôi đã nghĩ tới nhưng chưa làm được\\
\bottomrule
\end{longtable}
\end{center}

Câu hỏi mở đầu thứ ba và thứ tư đều nằm trong bảng này. \emph{Further work is needed} hầu như luôn đánh dấu một chỗ chưa xong --- và nếu nó xuất hiện ngay sau phần bàn về một kết quả yếu, đó là cách nói ``chúng tôi không giải thích được''. Còn \emph{would benefit from} trong nhận xét phản biện: về hình thức là gợi ý, về thực tế là yêu cầu --- người phản biện đang lịch sự, nhưng nếu bạn không làm và không giải thích vì sao thì lần nộp sau sẽ gặp lại đúng nhận xét đó.

\begin{cambay}
Đừng đảo ngược bảng này khi \emph{viết}. Nhiều người học, sau khi nhận ra các quy ước, bắt đầu dùng chúng để nói vòng vo mọi thứ --- và kết quả là bài viết không nói được gì rõ ràng. Quy ước phê bình lịch sự dùng cho \emph{phê bình người khác}; còn khi trình bày kết quả của chính mình, hãy nói đúng mức: mạnh khi có bằng chứng mạnh, dè dặt khi bằng chứng yếu, và luôn cụ thể về phạm vi. Chương \ptr{C/14} nói kỹ cách hiệu chỉnh.
\end{cambay}

\section{Đọc những gì không được nói}

Câu hỏi mở đầu thứ năm hỏi cách nhận ra thông tin nằm ở chỗ vắng mặt. Bốn chỗ đáng kiểm tra một cách có hệ thống.

\emph{Chỉ số vắng mặt.} Nếu một bài về hệ thống thời gian thực không báo cáo thời gian chạy, hoặc một bài về độ chính xác không báo cáo phương sai giữa các lần chạy, hãy giả định rằng con số đó không thuận lợi. Cùng lập luận Grice: tác giả có động cơ nêu điều tốt.

\emph{Đường cơ sở vắng mặt.} Một phương pháp mạnh và nổi tiếng mà không xuất hiện trong bảng so sánh, không kèm lý do, là dấu hiệu rõ.

\emph{Điều kiện vắng mặt.} Bài báo cáo kết quả nhưng không nói dữ liệu được thu thập thế nào, hoặc không nói tham số được chọn ra sao. Chỗ vắng mặt này thường ở phần phương pháp và nó ảnh hưởng trực tiếp tới tính tái lập (\ptr{B/07}).

\emph{Trường hợp thất bại vắng mặt.} Bài không có ví dụ nào về khi phương pháp không hoạt động. Với người đọc có kinh nghiệm, điều này làm giảm --- chứ không tăng --- độ tin cậy.

Đối lập với vắng mặt là các cách nói \emph{gần như phủ định} mà người học dễ đọc nhầm. \emph{We did not observe any degradation} không giống \emph{there is no degradation}: câu đầu nói về những gì họ quan sát được trong thí nghiệm của họ, câu sau là một khẳng định phổ quát. Tương tự, \emph{no significant difference was found} nghĩa là ``chúng tôi không tìm thấy khác biệt đủ lớn để kết luận'' chứ không phải ``hai bên như nhau'' --- một trong những hiểu nhầm phổ biến nhất khi đọc kết quả thống kê.

\section{Đọc nhận xét phản biện}

Đây là ứng dụng thực dụng nhất của chương và là việc bạn sẽ làm nhiều nếu công bố nghiên cứu.

Ba mức nghiêm trọng, đọc qua ngôn ngữ. \emph{Nhận xét bắt buộc:} dùng động từ mạnh hoặc nói về tính đúng đắn --- \emph{The claim in Section 4 is not supported by}, \emph{The comparison is unfair because}. Phải xử lý, hoặc phải phản biện lại có căn cứ. \emph{Nhận xét nên làm:} \emph{would benefit from}, \emph{it would be helpful to}, \emph{I suggest}. Nên làm, và nếu không làm thì phải giải thích. \emph{Nhận xét tuỳ chọn:} \emph{minor}, \emph{typo}, \emph{consider rephrasing}. Sửa nhanh và không cần bàn.

Hai mẫu đáng đọc kỹ. Khi phản biện viết \emph{I may have missed this, but\ldots} --- đó thường là lịch sự, và ý thật là ``chỗ này viết không rõ nên tôi không tìm thấy''; kể cả khi thông tin thực sự có trong bài, việc người đọc chuyên môn không tìm thấy đã là một vấn đề cần sửa. Và khi phản biện viết \emph{the authors should clarify whether\ldots}, hãy đọc như một nghi ngờ về tính đúng đắn, không chỉ về cách diễn đạt.

\begin{cannho}
Ba câu hỏi để đọc mức cam kết của bất kỳ khẳng định nào: (i) động từ báo cáo nào được dùng --- \emph{show} hay \emph{suggest}; (ii) có động từ tình thái hoặc trạng từ dè dặt nào không; (iii) phạm vi được giới hạn tới đâu --- có cụm \emph{in our experiments} hay không. Ba câu này mất mười giây và cho bạn biết chính xác tác giả đang cam kết điều gì --- thông tin mà bản dịch máy và cách đọc lướt gần như luôn làm mất.
\end{cannho}

\section{Gốc rễ: từ triết học ngôn ngữ tới văn phong khoa học}

Ý tưởng rằng nghĩa của một phát ngôn vượt ra ngoài nghĩa đen của các từ được hệ thống hoá bởi Grice trong loạt bài giảng cuối thập niên 1960, công bố năm 1975\cite{grice1975}. Bối cảnh là một cuộc tranh luận trong triết học: các nhà lôgic học phàn nàn rằng ngôn ngữ tự nhiên ``không chính xác'' vì \emph{and} đôi khi hàm ý thứ tự thời gian, \emph{or} đôi khi loại trừ. Grice lật ngược vấn đề: ngôn ngữ tự nhiên không thiếu chính xác, mà người nghe \emph{cộng thêm} một tầng suy luận dựa trên giả định hợp tác. Nghĩa đen và nghĩa được truyền đạt là hai thứ khác nhau, và cả hai đều có quy luật.

Nhánh thứ hai đến từ nghiên cứu về lịch sự. Brown và Levinson mô tả cách các ngôn ngữ khác nhau đều phát triển những cơ chế để giảm mức đe doạ khi phải nói điều bất lợi cho người nghe\cite{brownlevinson1987}. Trong học thuật, việc chỉ ra sai sót của đồng nghiệp là một hành động đe doạ điển hình --- và toàn bộ bộ quy ước ở mục 3 là cơ chế giảm nhẹ nó.

Nhánh thứ ba là nghiên cứu trực tiếp về văn phong khoa học. Phân tích hàng trăm bài báo, Hyland cho thấy rằng dè dặt và nhấn mạnh không phải trang trí mà là bộ phận chức năng của lập luận khoa học\cite{hyland1998,hyland2005}: chúng vừa mô tả mức bất định, vừa định vị tác giả trong cộng đồng, vừa mời người đọc tham gia đánh giá thay vì áp đặt. Ông cũng chỉ ra rằng mật độ và kiểu dè dặt khác nhau đáng kể giữa các ngành --- các ngành thực nghiệm dè dặt nhiều hơn toán học, vì bản chất bằng chứng khác nhau.

Với người học, ba nhánh này gộp lại thành một kết luận thực dụng: các quy ước ở chương này \emph{không} phải sự vòng vo tuỳ tiện. Chúng mã hoá thông tin thật --- mức chắc chắn, phạm vi, và vị thế của tác giả --- và đọc được chúng là đọc được một tầng nội dung mà bản dịch từng chữ sẽ làm mất sạch.

\begin{gocre}
\gr{Cuối thập niên 1960, công bố 1975 --- Grice}{Các nhà lôgic cho rằng ngôn ngữ tự nhiên thiếu chính xác.}{Nguyên lý hợp tác và hàm ý: nghĩa được truyền đạt khác nghĩa đen, và phần chênh có quy luật.}
\gr{1978, mở rộng 1987 --- Brown \& Levinson}{Vì sao mọi ngôn ngữ đều có cách nói vòng khi phải nói điều bất lợi?}{Lý thuyết lịch sự: các chiến lược giảm nhẹ hành động đe doạ thể diện --- nền của quy ước phê bình học thuật.}
\gr{Thập niên 1990--2000 --- Hyland}{Văn khoa học được coi là khách quan và vô ngã.}{Dè dặt và nhấn mạnh là bộ phận chức năng của lập luận, không phải trang trí; mật độ khác nhau theo ngành.}
\gr{Thực tiễn bình duyệt hiện đại}{Phê bình trực tiếp gây xung đột và làm hỏng quá trình đánh giá.}{Bộ quy ước diễn đạt phê bình gián tiếp --- hiệu quả trong cộng đồng quen thuộc, nhưng là rào cản với người mới và người không dùng tiếng Anh bản ngữ.}
\end{gocre}

\section{Liên hệ ngang}

\begin{lienhe}
\lh{Tiếng Việt}{Nói giảm nói tránh, ``có lẽ'', ``e rằng''}{Cùng cơ chế lịch sự nhưng \emph{phân bố khác}: tiếng Việt dùng nhiều tiểu từ tình thái và cách xưng hô, tiếng Anh học thuật dùng động từ tình thái và cấu trúc bị động. Người Việt thường dè dặt \emph{quá mức} khi viết tiếng Anh vì chuyển thẳng thói quen sang.}
\lh{Ngoại giao}{Ngôn ngữ công hàm: \emph{we note with concern}}{Cùng bộ mã hoá mức độ: từ \emph{note} tới \emph{express concern} tới \emph{condemn} là một thang rõ ràng mà giới trong nghề đọc chính xác. Cùng rủi ro: người ngoài đọc theo nghĩa đen sẽ hiểu sai mức nghiêm trọng.}
\lh{Đánh giá mã nguồn}{\emph{Nit:}, \emph{Consider\ldots}, \emph{This will break when\ldots}}{Cùng ba mức nghiêm trọng như ở mục 6. Nhiều nhóm nay quy ước tiền tố rõ ràng (\emph{nit}, \emph{question}, \emph{blocking}) chính để tránh việc người đọc phải giải mã --- một cải tiến mà học thuật chưa có.}
\lh{Phản hồi trong công việc}{Đánh giá hiệu suất, thư giới thiệu}{Thư giới thiệu tiếng Anh có quy ước đọc riêng: điều \emph{không} được khen mới là thông tin. Cùng nguyên lý Grice ở mục 1.}
\lh{Y khoa}{Thông báo tin xấu, ghi chép bệnh án}{Ngôn ngữ được hiệu chỉnh cẩn thận về mức chắc chắn: \emph{consistent with}, \emph{cannot be excluded}. Cùng chức năng: mô tả đúng mức bất định và giới hạn trách nhiệm.}
\lh{Pháp lý}{Ngôn ngữ hợp đồng và bản án}{Ngược lại với học thuật: pháp lý cố loại bỏ hàm ý bằng cách định nghĩa mọi thứ. So sánh hai thể loại cho thấy rõ chức năng của dè dặt --- ở đâu bất định là thật thì phải nói ra, ở đâu không được phép bất định thì phải định nghĩa.}
\end{lienhe}

\begin{traloi}
\tl{Bạn suy ra rằng độ chính xác không phải điểm mạnh. Nguyên tắc là nguyên lý hợp tác của Grice: trong tóm tắt, tác giả có mọi động cơ nêu ưu điểm và có đủ chỗ để nêu, nên việc \emph{không} nhắc một khía cạnh quan trọng cho phép người đọc suy ra rằng khía cạnh đó không thuận lợi. Câu văn không sai sự thật; thông tin nằm ở lựa chọn nói gì và im lặng ở đâu.}
\tl{Ba lý do đều chính đáng. Nhận thức: kết quả thực nghiệm luôn có phạm vi hữu hạn, nên dè dặt là mô tả đúng mức bất định. Xã hội: khẳng định mạnh mà sai làm hỏng uy tín; khẳng định hiệu chỉnh đúng thì vẫn đứng vững. Tương tác: dè dặt để chỗ cho người đọc không đồng ý mà không phải bác bỏ tác giả. Điều cần phân biệt khi đọc: dè dặt \emph{cụ thể} (\emph{in our experiments}) mang thông tin thật, còn dè dặt nghi thức (\emph{it may be argued that}) chỉ là quy ước thể loại.}
\tl{Nó đánh dấu một vấn đề chưa được giải quyết, và tuỳ vị trí mà mức nghiêm trọng khác nhau. Ở cuối phần kết luận, nó thường là nghi thức. Nhưng ngay sau phần bàn về một kết quả yếu hoặc một hiện tượng bất thường, nó thường có nghĩa ``chúng tôi không giải thích được điều này''. Cách kiểm: xem câu trước nó nói gì --- \emph{further work} nối vào một kết quả bất lợi là dấu hiệu đáng chú ý hơn nhiều so với \emph{further work} nối vào một hướng mở rộng chung chung.}
\tl{Về hình thức là gợi ý, về thực tế là yêu cầu. Trong bình duyệt, đây là cách diễn đạt lịch sự chuẩn cho ``hãy làm việc này''. Nếu bạn không làm, bạn phải giải thích rõ vì sao trong thư phản hồi --- và lý do phải có căn cứ, không chỉ là ``ngoài phạm vi''. Ba mức nghiêm trọng ở mục 6: \emph{is not supported by} là bắt buộc; \emph{would benefit from} là nên làm; \emph{consider rephrasing} là tuỳ chọn.}
\tl{Kiểm bốn chỗ vắng mặt một cách có hệ thống: chỉ số lẽ ra phải báo cáo mà không có (thời gian chạy trong bài về hệ thống thời gian thực, phương sai giữa các lần chạy); đường cơ sở mạnh không xuất hiện mà không kèm lý do; điều kiện thí nghiệm không được mô tả đủ để tái lập; và không có ví dụ nào về trường hợp thất bại. Ngoài ra, chú ý phân biệt \emph{we did not observe X} với \emph{there is no X}, và hiểu rằng \emph{no significant difference} nghĩa là không đủ bằng chứng để kết luận khác biệt, không phải hai bên như nhau.}
\end{traloi}

\begin{dongian}
Trong tiếng Anh học thuật, người ta hiếm khi nói thẳng điều tiêu cực. Không phải vì họ vòng vo, mà vì cộng đồng đã thống nhất một bộ mã --- và người trong nghề đọc bộ mã đó chính xác như đọc chữ.

Ví dụ dễ nhất: \emph{Smith và cộng sự báo cáo kết quả ấn tượng trên các chuỗi dữ liệu trong nhà}. Nghe như khen. Nhưng chữ ``trong nhà'' làm việc của nó: nếu phương pháp kia cũng tốt ngoài trời thì đã không cần giới hạn như vậy. Người đọc quen sẽ hiểu ngay câu tiếp theo sắp nói ``còn ngoài trời thì chưa ai làm''.

Nguyên tắc chung rất đơn giản: \emph{nếu người viết có thể nói điều mạnh hơn mà không nói, tức là họ nói không được}. Vì thế một bài khoe ``phương pháp của chúng tôi nhanh'' mà không nhắc gì tới độ chính xác thì gần như chắc chắn là độ chính xác không nổi bật.

Cũng vì vậy mà những từ nhỏ như \emph{may}, \emph{suggest}, \emph{in our experiments} không phải là văn phong rườm rà. Chúng nói cho bạn biết tác giả tin điều này tới mức nào và trong phạm vi nào. Một câu viết \emph{these results suggest} khác hẳn \emph{these results show}, và bản dịch tiếng Việt thường xoá mất khác biệt đó.

Cuối cùng, một bảng nhỏ đáng thuộc khi đọc nhận xét của người phản biện. \emph{The paper would benefit from X} nghĩa là ``hãy làm X''. \emph{It is not entirely clear how\ldots} nghĩa là ``chỗ này bạn viết không rõ, hoặc lập luận có lỗ hổng''. \emph{Further work is needed} nghĩa là ``chưa xong''. Và \emph{I may have missed this, but\ldots} gần như luôn nghĩa là ``bạn viết không rõ nên tôi không tìm thấy'' --- và kể cả khi thông tin có thật trong bài, việc người đọc chuyên môn không tìm thấy đã là vấn đề của bài.
\motcau{Nếu tác giả có thể nói mạnh hơn mà không nói, sự im lặng đó chính là thông tin.}
\chuadongian{Các quy ước ở đây thay đổi theo ngành, theo tạp chí và theo cá nhân; chúng là xác suất cao chứ không phải quy tắc chắc chắn, nên đừng dùng chúng để kết tội ai.}
\end{dongian}

\begin{onbai}
\ar[3]{Nguyên lý hợp tác của Grice}{Người nói được giả định là nói đủ, đúng, liên quan, rõ. Vì thế điều \emph{không} được nói mang thông tin: nói mạnh hơn được mà không nói tức là không nói được.}
\ar[3]{Ba câu hỏi đọc mức cam kết}{Động từ báo cáo nào; có tình thái hoặc trạng từ dè dặt không; phạm vi bị giới hạn tới đâu.}
\ar[3]{Ba lý do khoa học dè dặt}{Nhận thức (mô tả đúng bất định); xã hội (bảo vệ uy tín); tương tác (để chỗ cho người đọc không đồng ý).}
\ar[3]{Dè dặt cụ thể so với nghi thức}{\emph{In our experiments} mang thông tin thật; \emph{it may be argued that} chỉ là quy ước. Phân biệt bằng tính cụ thể.}
\ar[3]{Bảng giải mã hay dùng}{\emph{Little attention} = chưa ai làm; \emph{further work is needed} = chưa xong; \emph{would benefit from} = hãy làm; \emph{competitive with} = chỉ ngang bằng.}
\ar[3]{Bốn chỗ vắng mặt cần kiểm}{Chỉ số lẽ ra phải có; đường cơ sở mạnh không xuất hiện; điều kiện thí nghiệm không đủ tái lập; không có trường hợp thất bại.}
\ar[3]{\emph{No significant difference}}{Nghĩa là không đủ bằng chứng để kết luận có khác biệt --- \emph{không} phải hai bên như nhau.}
\ar[2]{Ba mức nhận xét phản biện}{Bắt buộc (\emph{not supported by}, \emph{unfair}); nên làm (\emph{would benefit from}, \emph{I suggest}); tuỳ chọn (\emph{minor}, \emph{typo}).}
\ar[2]{\emph{I may have missed this, but}}{Lịch sự cho ``chỗ này viết không rõ''; kể cả khi thông tin có trong bài, việc người đọc chuyên môn không thấy đã là vấn đề.}
\ar[2]{Nói giảm}{\emph{Not uncommon} = khá phổ biến; \emph{not unreasonable} = hợp lý. Thường mạnh hơn nghĩa đen, phổ biến trong văn Anh.}
\ar[2]{\emph{Clearly} và \emph{obviously}}{Đôi khi thay cho lập luận còn thiếu --- gặp chúng thì nên hỏi ``chỗ này có thật sự hiển nhiên không''.}
\ar[2]{Cạm bẫy khi viết}{Đừng dùng quy ước phê bình lịch sự cho \emph{kết quả của chính mình}; trình bày kết quả cần nói đúng mức và cụ thể về phạm vi (\ptr{C/14}).}
\ar[1]{Khác biệt theo ngành}{Mật độ dè dặt khác nhau: ngành thực nghiệm dè dặt nhiều hơn toán học, vì bản chất bằng chứng khác nhau.}
\end{onbai}

\begin{hoidap}
\qa{Làm sao biết một dè dặt là thật hay chỉ là nghi thức?}{Kiểm tính cụ thể và vị trí. Dè dặt thật thường \emph{cụ thể} và gắn với một điều kiện kiểm chứng được: \emph{on the three datasets considered}, \emph{when the baseline is tuned per sequence}. Dè dặt nghi thức thì chung chung và có thể bỏ đi mà câu không mất thông tin: \emph{it may be argued that}, \emph{to some extent}. Ngoài ra hãy xem vị trí: dè dặt trong phần kết quả thường có nội dung; dè dặt trong phần kết luận thường là nghi thức.}
\qa{Người Việt viết tiếng Anh thường dè dặt quá mức. Vì sao và sửa thế nào?}{Vì thói quen lịch sự trong tiếng Việt được chuyển thẳng sang, cộng với tâm lý e ngại khi viết bằng ngoại ngữ. Kết quả là những câu như \emph{we think that maybe our method could possibly improve} --- chồng ba lớp dè dặt, và người đọc hiểu là bạn không tin vào kết quả của chính mình. Cách sửa: mỗi khẳng định chỉ dùng \emph{một} lớp dè dặt; và thay dè dặt mơ hồ bằng giới hạn phạm vi cụ thể --- \emph{on our test set, the method reduces error by 12\%} mạnh và trung thực hơn \emph{our method may possibly be better}.}
\qa{Khi phản biện bài người khác, tôi nên dùng các mẫu lịch sự này tới mức nào?}{Dùng đủ để không xúc phạm, nhưng đừng để mất tính rõ ràng. Cấu trúc an toàn cho một nhận xét nghiêm trọng gồm ba phần: nêu điều bạn quan sát được một cách trung tính (\emph{The comparison in Table 2 uses default parameters for the baseline}); nói hệ quả (\emph{This may overstate the improvement}); và đề xuất cụ thể (\emph{Tuning the baseline on the validation set would make the comparison fairer}). Ba phần này lịch sự mà không mơ hồ --- và người nhận biết chính xác phải làm gì.}
\qa{Có phải mọi cách nói vòng đều mang hàm ý không?}{Không. Nhiều cụm chỉ là quy ước thể loại đã mòn nghĩa --- \emph{it should be noted that}, \emph{as mentioned above} --- và cố tìm hàm ý ở đó là đọc quá mức. Dấu hiệu để phân biệt: hàm ý thật thường xuất hiện ở những chỗ \emph{có gì đó để mất} --- khi đánh giá công trình người khác, khi báo cáo kết quả không thuận lợi, khi thừa nhận hạn chế. Ở phần mô tả phương pháp thuần tuý, ngôn ngữ thường nói đúng nghĩa đen.}
\qa{Tôi đọc bản dịch máy của bài báo có mất gì không?}{Mất chính tầng của chương này. Các công cụ dịch xử lý nghĩa quy chiếu khá tốt nhưng thường san phẳng mức cam kết: \emph{suggest} và \emph{show} có thể ra cùng một từ tiếng Việt, các cụm giới hạn phạm vi bị lược, và nói giảm bị dịch theo nghĩa đen. Nếu phải dùng bản dịch để nắm nội dung nhanh, hãy quay lại bản gốc cho ba chỗ: tóm tắt, các câu kết luận chính, và phần hạn chế --- đó là ba chỗ mật độ dè dặt cao nhất.}
\qa{Làm sao luyện kỹ năng này?}{Một bài tập rất hiệu quả và làm được mỗi ngày: lấy một đoạn trong bài báo, gạch chân mọi phương tiện dè dặt và nhấn mạnh, rồi \emph{viết lại đoạn đó ở mức cam kết cao nhất} --- bỏ hết dè dặt. So hai bản và bạn sẽ thấy chính xác tác giả đã hiệu chỉnh những gì. Làm việc này với hai mươi đoạn trong hai tuần sẽ đổi hẳn cách bạn đọc, vì bạn bắt đầu nhìn thấy các phương tiện đó thay vì lướt qua chúng.}
\qa{Bảng giải mã ở mục 3 có đúng cho mọi ngành không?}{Không hoàn toàn. Mức độ gián tiếp thay đổi theo ngành và theo văn hoá học thuật: một số cộng đồng phê bình thẳng hơn nhiều, và các bài đánh giá công khai gần đây có xu hướng trực tiếp hơn so với bình duyệt kín truyền thống. Cách hiệu chỉnh bảng cho ngành của bạn: đọc mười phần công trình liên quan trong đúng tạp chí bạn nhắm tới và ghi lại cách họ diễn đạt phê bình --- sau mười bài, bạn sẽ có bảng riêng chính xác hơn bảng chung.}
\end{hoidap}

\begin{baitap}
\bt[1]{Gạch chân mọi phương tiện dè dặt và nhấn mạnh trong một phần tóm tắt}{Bài báo ngành}{Với mỗi cái, ghi nó làm mạnh hay làm nhẹ khẳng định.}
\bt[1]{Viết lại một đoạn ở mức cam kết cao nhất, bỏ hết dè dặt}{Bài tập viết}{So với bản gốc để thấy tác giả đã hiệu chỉnh những gì.}
\bt[2]{Tìm trong một bài ba chỗ có thông tin nằm ở sự vắng mặt}{Tài liệu ngành}{Chỉ số thiếu, đường cơ sở thiếu, hoặc trường hợp thất bại không được báo cáo.}
\bt[2]{Lập bảng giải mã riêng cho ngành của bạn từ 10 phần công trình liên quan}{Tài liệu ngành}{Ghi lại các mẫu phê bình lịch sự thật sự được dùng trong tạp chí bạn nhắm tới.}
\bt[2]{Phân loại các nhận xét trong một bản phản biện thành ba mức nghiêm trọng}{Phản biện bạn từng nhận}{Nếu chưa có, tìm các bản phản biện công khai; nhiều hội nghị nay công bố chúng.}
\bt[3]{Viết ba nhận xét phản biện theo cấu trúc quan sát--hệ quả--đề xuất}{Bài tập viết}{Cho một bài bạn vừa đọc; mục tiêu là lịch sự mà vẫn rõ ràng.}
\bt[3]{Soát lại một đoạn bạn từng viết và đếm số lớp dè dặt trong mỗi câu}{Bài viết của bạn}{Sửa để mỗi khẳng định chỉ có một lớp, và thay dè dặt mơ hồ bằng giới hạn phạm vi cụ thể.}
\end{baitap}

\section*{Kết chương và đọc tiếp}

Đây là chương gần nhất với mục tiêu ``hiểu đúng và sâu ý người nói''. Ba điều đáng mang theo: nguyên lý ``im lặng là thông tin'' của Grice; ba câu hỏi mười giây để định vị mức cam kết của một khẳng định; và bảng giải mã phê bình lịch sự, kèm cảnh báo rằng nó là quy ước xác suất cao chứ không phải quy tắc chắc chắn.

Hai chương còn lại của phần B áp bộ công cụ này vào hai tình huống cụ thể. Chương \ptr{B/09} nói về nghe giảng và xem khoá học --- nơi bạn không dừng lại tra được và phải bắt cấu trúc theo thời gian thực. Chương \ptr{B/10} khép lại phần đọc bằng việc biến những gì bạn đọc thành vốn từ chuyên ngành của riêng bạn.
\end{document}
